\section{\sname: A Coupled Lifelong Fleet Manager}\label{sec:method}

\sname\ runs a rolling-horizon loop (\cref{fig:overview}): every replanning
epoch it (i) re-plans the robots that already have goals, (ii) sends
low-charge robots to free charger slots, (iii) prices every idle robot against
the pending tasks and solves an assignment problem, and (iv) plans the newly
assigned robots into the same reservation table. The coupling lives in step
(iii): assignment costs are not free-space distances but arrival times
computed by the same windowed planner against the reservations written in
step (i). We first fix the model (\cref{subsec:model}), then describe the
planning layer (\cref{subsec:planning}), the charging policy
(\cref{subsec:charging}) and the coupled allocation (\cref{subsec:alloc}).

\subsection{Modelling the fleet: lifelong MAPD with battery constraints}\label{subsec:model}

\noindent\textbf{Workspace.}
The workspace is a 4-connected grid graph $G=(V,E)$ with obstacle cells
removed. Four endpoint classes mirror the operator console's vocabulary:
pickup points $\mathcal{P}\subset V$ on shelf lines, delivery docks
$\mathcal{D}\subset V$, charger slots $\mathcal{C}\subset V$ grouped into
stations, and home bays $\mathcal{H}\subset V$, one per robot. Every dock,
slot and bay is a dead-end pocket off a corridor, so an instance is
well-formed in the sense of \citet{ma2017lifelong}: a robot resting on a
non-task endpoint never blocks a path between two other endpoints.
Distances $d(u,v)$ are exact shortest-path lengths on $G$, computed by
breadth-first search from each endpoint and cached.

\noindent\textbf{Robots and tasks.}
A fleet $\mathcal{R}=\{1,\dots,n\}$ of robots moves synchronously in discrete
time $t\in\mathbb{N}$; at every tick a robot either moves to a 4-neighbour or
waits. Robot $r$ has position $p_t^r\in V$, state of charge
$s_t^r\in[0,1]$, and home bay $h^r\in\mathcal{H}$. Tasks arrive as a stream
$\mathcal{T}=\{\tau_k\}$, $\tau_k=(a_k, p_k, d_k)$ with arrival time $a_k$,
pickup $p_k\in\mathcal{P}$ and dock $d_k\in\mathcal{D}$; a robot that reaches
$p_k$ loads for $\delta$ ticks, drives to $d_k$ and unloads for $\delta$
ticks, which completes the task at time $c_k$.

\noindent\textbf{Battery.}
Charge decreases by $\beta_{\mathrm{m}}$ per move and $\beta_{\mathrm{w}}$ per
wait, and increases by $\gamma$ per tick while a robot rests on a charger
slot; a robot with $s^r_t\le 0$ is depleted and becomes a permanent
obstacle. The charging policy is a two-threshold hysteresis
$(\theta_{\mathrm{lo}},\theta_{\mathrm{hi}})$, \cref{subsec:charging}.

\noindent\textbf{Conflicts.}
Two robots conflict when they occupy the same cell at the same tick (vertex
conflict) or traverse the same edge in opposite directions during the same
tick (swap conflict); following a robot into the cell it just left is
allowed, as in the standard MAPF definitions of \citet{stern2019mapf}.

\noindent\textbf{Objective.}
Given the stream $\mathcal{T}$, the fleet manager chooses, at every tick, an
assignment of idle robots to pending tasks, a charging decision and one action
per robot, so as to minimise the mean service time of the delivered tasks
\begin{equation}\label{eq:objective}
  J \;=\; \frac{1}{|\mathcal{T}|}\sum_{\tau_k\in\mathcal{T}} \bigl(c_k - a_k\bigr)
  \quad\text{subject to}\quad
  \text{no conflicts},\;\; s^r_t>0\;\;\forall r,t,
\end{equation}
and we also report throughput $|\mathcal{T}|/T_{\mathrm{mk}}$, where the
makespan $T_{\mathrm{mk}}=\max_k c_k$. For a finite stream throughput is the
inverse makespan, a different objective from mean service time that can rank
schedules differently, so we report both. Optimising \cref{eq:objective} exactly is intractable even
without batteries \citep{yu2013structure}; \sname\ is a rolling-horizon
heuristic. It enforces the conflict constraint by construction
(\cref{prop:pp,cor:exec}); the charge constraint is only guarded by a
heuristic feasibility test (\cref{eq:feasible}), and progress (absence of
deadlock and livelock) is not guaranteed. Both are measured in
\cref{sec:exp}.

\subsection{Planning under reservations: windowed space-time A* with a hold cascade}\label{subsec:planning}

\noindent\textbf{Reservation table.}
A reservation table $\mathcal{B}$ stores, for an epoch starting at $t_0$ with
horizon $t_0+w$, the vertex reservations $(v,t)$ and edge reservations
$(u,v,t)$ of every planned robot, plus \emph{holds} $(v,t_{\mathrm{h}})$
meaning ``occupied from $t_{\mathrm{h}}$ until the horizon''. A path whose
last cell is reached at $t_{\mathrm{e}}\le t_0+w$ ends with a hold, so a
robot that arrives early is a static obstacle for the rest of the window.
The table answers three queries: $\mathrm{free}(v,t)$,
$\mathrm{free}(u\!\to\!v,t)$ (false iff the reverse edge is reserved at $t$)
and $\mathrm{goalfree}(v,t)$ (true iff $v$ is free for all
$t'\in[t,t_0+w]$ and not held).

\noindent\textbf{Windowed space-time A*.}
Each robot plans on the time-expanded graph
$V\times\{t_0,\dots,t_0+w\}$ with the exact distance $d(\cdot,g)$ as
heuristic (\cref{alg:astar}). A state $(g,t)$ is a goal only if
$\mathrm{goalfree}(g,t)$; a state at the horizon is closed with cost
$t-t_0+d(v,g)$ and completed with a deterministic shortest path, because no
reservation exists beyond the horizon. The returned cost is the arrival time
\begin{equation}\label{eq:arrival}
  A^{\mathcal{B}}(p,g,t_0) \;=\; \min\{\,t - t_0 : (g,t)\ \text{reachable under }\mathcal{B},\ \mathrm{goalfree}(g,t)\,\},
\end{equation}
the per-agent cost of sum-of-costs MAPF restricted to the window. Since
$d(\cdot,g)$ is consistent, the first goal state popped is optimal, and the
search is bounded by $|V|(w+1)$ expansions.

\begin{algorithm}[t]
\caption{Windowed space-time A* (one robot)}\label{alg:astar}
\KwIn{start $p$, goal $g$, epoch start $t_0$, window $w$, reservations $\mathcal{B}$, distances $d(\cdot,g)$}
\KwOut{path $\pi=(\pi_0=p,\pi_1,\dots)$ and arrival time, or \textsc{fail}}
$\mathit{open}\gets\{(p,t_0)\}$ with $f=d(p,g)$; $\mathit{closed}\gets\emptyset$\;
\While{$\mathit{open}\neq\emptyset$}{
  pop $(v,t)$ with minimal $f$ (ties: larger $g$-value, then FIFO)\;
  \lIf{$(v,t)\in\mathit{closed}$}{\textbf{continue}}
  add $(v,t)$ to $\mathit{closed}$\;
  \lIf{$v=g$ \KwSty{and} $\mathrm{goalfree}(g,t)$}{\Return path to $(v,t)$, cost $t-t_0$}
  \lIf{$t=t_0+w$}{\Return path to $(v,t)$ followed by a shortest path $v\!\leadsto\! g$, cost $t-t_0+d(v,g)$}
  \ForEach{$v'\in\{v\}\cup N(v)$}{
    \If{$\mathrm{free}(v',t+1)$ \KwSty{and} $(v'=v$ \KwSty{or} $\mathrm{free}(v\!\to\!v',t))$}{
      push $(v',t+1)$ with $f=(t+1-t_0)+d(v',g)$\;
    }
  }
}
\Return \textsc{fail}\;
\end{algorithm}

\noindent\textbf{Prioritized planning with a hold cascade.}
Robots are planned one after another in priority order
(\cref{alg:pp}); each path is written into $\mathcal{B}$ before the next
robot plans. Prioritized planning is incomplete: a robot may find no path
within the window. Instead of leaving such a robot unplanned, \sname\ makes it
\emph{hold} its current cell for the whole window, and re-plans every robot
whose reservation crosses that cell---including robots planned in an earlier
group of the same epoch. Holds only accumulate within an epoch, so the cascade
terminates, and at termination every path is consistent with every
reservation:

\begin{proposition}[Conflict-freeness of an epoch]\label{prop:pp}
Let all robots occupy distinct cells at $t_0$, and let every entry of
$\mathcal{B}$ at the call belong either to an external robot in
$\mathcal{X}$ or to a static robot holding its own current cell. If the
entries of $\mathcal{B}$ are pairwise conflict-free when \cref{alg:pp} is
called, then after it returns the set of paths in $\mathcal{B}$ is vertex-
and swap-conflict-free on $[t_0,t_0+w]$, and \cref{alg:pp} calls the
low-level planner at most $N(N+1)/2\le|\mathcal{R}|(|\mathcal{R}|+1)/2$
times, where $N=m+|\mathcal{X}|$.
\end{proposition}
\begin{proof}
\emph{Invariant.} We show that no two entries of $\mathcal{B}$ conflict at
any moment. A path is inserted only when \cref{alg:astar} returned it under
the current $\mathcal{B}$: every vertex $(\pi_k,t_0+k)$ passed
$\mathrm{free}$, every move passed the swap query, and, if the path reaches
its goal inside the window, the hold at its last cell passed
$\mathrm{goalfree}$, so no existing entry occupies that cell later in the
window (a path that reaches the horizon first is stored only up to the
horizon, without a hold, and its unconstrained tail is never executed); entries inserted afterwards respect the hold because
$\mathrm{free}$ checks holds. A hold for a failed robot $r$ is written at
$p^r$ from $t_0$ on. No static robot and no held robot occupies $p^r$ (their
cells are their own, and start cells are distinct), and by assumption every
other entry belongs to a robot in $\mathcal{X}$ or in the input list; all of
those that visit $p^r$ at some $t\in[t_0,t_0+w]$ form $\mathcal{V}$ and are
removed before the hold is written. A hold never moves, so it cannot take
part in a swap conflict. Removing entries cannot create a conflict, so the
invariant holds at termination.

\emph{Call bound.} A robot is popped once initially if it is in the input
list ($m$ calls) and once each time it is re-queued. Re-queuing happens only
when a hold is created; a held robot is never re-queued (it is excluded from
$\mathcal{V}$) and so fails at most once, hence at most $N$ holds are
created. When the $k$-th hold is created, $|\mathcal{K}|=k$ and
$\mathcal{V}$ contains only robots outside $\mathcal{K}$, so
$|\mathcal{V}|\le N-k$. The number of calls is therefore at most
$m+\sum_{k=1}^{N}(N-k)\le N+N(N-1)/2=N(N+1)/2$.
\end{proof}

The assumption on $\mathcal{B}$ is what the epoch
(\cref{alg:pact}) maintains: static robots are written as holds on their own
cells, group~1 is planned into a table that contains only those holds, and
group~2 is planned with group~1 as $\mathcal{X}$. The optional CBS backend
(below) only inserts paths that are conflict-free with each other and with
$\mathcal{B}$ inside the window, which preserves the invariant, and the
robots it wrote into $\mathcal{B}$ are added to $\mathcal{X}$ of the
subsequent call of \cref{alg:pp} for the rest of the group. Neither the proposition nor its corollary says anything about
progress: a robot may hold for several epochs.

\begin{corollary}[Conflict-free execution]\label{cor:exec}
If $h\le w$, every step executed by the simulator is vertex- and
swap-conflict-free.
\end{corollary}
\begin{proof}
Every robot executes the path it has in the current table: a planned path,
or a hold on its own cell for static robots, and a robot whose path has ended
rests on its last cell, which that path holds. Epochs start at least every
$h$ ticks, so the step from $t$ to $t+1$ executed after an epoch at $t_0$
satisfies $t+1\le t_0+h\le t_0+w$ and lies inside the window in which the
table is conflict-free by \cref{prop:pp}. Between epochs the table changes
only by single-robot re-plans (below): the robot's old entry is removed, and
the new one is either a path returned by \cref{alg:astar} under the table or,
if that fails, a hold on the cell the robot stands on, which its old entry
held from the current tick on, so no other entry uses it. A depleted robot
stops, and an epoch starts before the next step is executed.
\end{proof}

\begin{algorithm}[t]
\caption{Prioritized planning with hold cascade}\label{alg:pp}
\KwIn{ordered robots $(r_1,\dots,r_m)$ with starts and goals, external robots $\mathcal{X}$ already in $\mathcal{B}$, epoch $t_0$}
\KwOut{paths for $r_1,\dots,r_m$ (and re-planned members of $\mathcal{X}$), held set $\mathcal{K}$}
$Q\gets(r_1,\dots,r_m)$; $\mathcal{K}\gets\emptyset$\;
\While{$Q\neq\emptyset$}{
  $r\gets$ pop front of $Q$\;
  $\pi\gets$ \textsc{A*}$(p^r,g^r,t_0,w,\mathcal{B})$ \tcp*{\cref{alg:astar}}
  \If{$\pi=$ \textsc{fail}}{
    $\pi\gets(p^r)$; $\mathcal{K}\gets\mathcal{K}\cup\{r\}$ \tcp*{hold current cell}
    $\mathcal{V}\gets\{\,r'\notin\mathcal{K} : r'$ reserved $(p^r,t)$ in $\mathcal{B}$ for some $t\in[t_0,t_0+w]\,\}$\;
    remove the reservations of every $r'\in\mathcal{V}$ from $\mathcal{B}$\;
    push $\mathcal{V}$ (in priority order) to the front of $Q$\;
  }
  write $\pi$ into $\mathcal{B}$ (with a hold at its last cell if it ends before the horizon)\;
}
\end{algorithm}

\noindent\textbf{Priorities and groups.}
An epoch plans two groups. Group~1 contains the busy robots (driving to a
pickup, a dock or a charger) together with parking robots that either stand
on a busy robot's goal or failed to plan in the previous epoch; group~2
contains the robots assigned in this epoch and the remaining parking robots.
Within a group, robots standing on another robot's goal come first, then
robots held last epoch, then carrying, then collecting, then charging-bound,
then parking robots, ties by id. Planning a held robot first is what evacuates
a robot boxed in on the single access cell of a corner dock; without it the
queue for that dock and the evacuating robot wait for each other
(\cref{sec:disc}). Robots that are charging, loading, unloading, depleted or
idle at home are static and are written as holds before group~1.

\noindent\textbf{Windowed CBS.}
As an optional backend, a group is planned jointly by Conflict-Based Search
\citep{sharon2015cbs} whose low level is \cref{alg:astar} under the union of
$\mathcal{B}$ and the node's constraints, and whose high level only resolves
conflicts inside the window (\cref{alg:cbs}). Two robots that share a goal
are not a valid one-shot instance; the lower-priority one is deferred to
\cref{alg:pp}, as are robots without an individually feasible path and the
whole group when the constraint tree exceeds a node budget.

\begin{algorithm}[t]
\caption{Windowed CBS for one group (high level)}\label{alg:cbs}
\KwIn{robots $\mathcal{A}$ with distinct goals, base reservations $\mathcal{B}$, epoch $t_0$, window $w$, node budget $N_{\max}$}
\KwOut{conflict-free paths on $[t_0,t_0+w]$ with minimal sum of arrival times, or \textsc{fail}}
root: $\mathcal{C}_r\gets\emptyset$, $\pi_r\gets$ \textsc{A*} under $\mathcal{B}\cup\mathcal{C}_r$ for each $r\in\mathcal{A}$\;
\lIf{some $\pi_r=$ \textsc{fail}}{\Return \textsc{fail}}
$\mathit{open}\gets\{\text{root}\}$\;
\While{$\mathit{open}\neq\emptyset$}{
  $\nu\gets$ node with minimal $\sum_r\mathrm{cost}(\pi_r)$ (ties: fewer conflicts, then FIFO)\;
  $\kappa\gets$ earliest conflict of $\nu$ on $[t_0,t_0+w]$\;
  \lIf{$\kappa=\emptyset$}{\Return paths of $\nu$}
  \lIf{$|\text{generated}|\ge N_{\max}$}{\Return \textsc{fail}}
  \ForEach{robot $r\in\kappa$}{
    child $\nu'\gets\nu$ with $\mathcal{C}_r$ extended by the vertex $(v,t)$ or edge $(u\!\to\!v,t)$ of $\kappa$\;
    re-plan $\pi_r$ under $\mathcal{B}\cup\mathcal{C}_r$; \lIf{feasible}{add $\nu'$ to $\mathit{open}$}
  }
}
\Return \textsc{fail}\;
\end{algorithm}

\noindent\textbf{Execution.}
Every tick each robot executes the next action of its path; a robot that
reaches the end of its path at its goal (and only then---a plan may pass
through the goal earlier while the goal is temporarily held) transitions
(pickup $\to$ loading, dock $\to$ unloading, slot $\to$ charging, bay $\to$
idle) and is immediately re-planned alone against the live table when it
gets a new goal; this re-plan keeps the table conflict-free
(\cref{cor:exec}). A new epoch starts every $h\le w$ ticks and whenever a
robot becomes idle while tasks are pending or a robot depletes. Independently
of the reservation table, the simulator compares the cells of all robots
before and after every executed step and counts vertex and swap conflicts;
this check tests the implementation against \cref{cor:exec}, and over all runs
in \cref{sec:exp} the count is zero.

\subsection{Charging: threshold policy with capacity-constrained slot assignment}\label{subsec:charging}

A robot with $s^r_t<\theta_{\mathrm{lo}}$ that is idle or parking must charge
before it accepts work; it leaves the slot when
$s^r_t\ge\theta_{\mathrm{hi}}$. Free slots are those neither occupied nor
already claimed by a robot en route, so a station with $k$ slots serves at
most $k$ robots at a time. Robots needing charge and free slots are matched
by the Hungarian method on $d(p^r,\sigma)$ (\cref{alg:charge}); a robot for
which no slot is free waits at its bay and is retried next epoch. A task
feasibility test guards against stranding: robot $r$ may take task $\tau$
only if
\begin{equation}\label{eq:feasible}
  s^r_t \;\ge\; \beta_{\mathrm{m}}\bigl(\hat a(r,\tau) + d(p_\tau,d_\tau) + \min_{\sigma\in\mathcal{C}} d(d_\tau,\sigma)\bigr) + \rho,
\end{equation}
where $\hat a$ is the estimated time to the pickup used by the allocation
(\cref{subsec:alloc}) and $\rho$ a reserve margin. The test is a heuristic,
not a guarantee: it budgets free-space moves only and leaves waits, detours,
dwell and queueing for a charger slot to the margin $\rho$, and a robot that
depletes nevertheless becomes a permanent obstacle. With
\cref{eq:feasible} and $\theta_{\mathrm{lo}}$ above the energy of the longest
bay-to-slot trip, no depletion occurred in any run of \cref{sec:exp}.

\begin{algorithm}[t]
\caption{Charger assignment (one epoch)}\label{alg:charge}
\KwIn{robots $\mathcal{N}=\{r : s^r_t<\theta_{\mathrm{lo}},\ r$ idle or parking$\}$, free slots $\mathcal{F}\subseteq\mathcal{C}$}
\lIf{$\mathcal{N}=\emptyset$ \KwSty{or} $\mathcal{F}=\emptyset$}{\Return}
$M\gets$ \textsc{Hungarian}$\bigl([\,d(p^r,\sigma)\,]_{r\in\mathcal{N},\sigma\in\mathcal{F}}\bigr)$\;
\ForEach{$(r,\sigma)\in M$}{claim $\sigma$ for $r$; goal of $r\gets\sigma$; status of $r\gets$ charging-bound\;}
\end{algorithm}

\subsection{Allocating tasks: assignment costs from MAPF-in-the-loop}\label{subsec:alloc}

\noindent\textbf{Candidate pool.}
Pending tasks are kept in arrival order. Following token passing
\citep{ma2017lifelong}, a task enters the candidate pool only while its dock
has fewer than $\kappa_{\mathrm{d}}$ robots in flight (assigned and not yet
delivered) and its pickup fewer than $\kappa_{\mathrm{p}}$; the pool holds
the oldest $\max(2m,8)$ such tasks for $m$ eligible robots. This endpoint
reservation is what keeps dock queues short enough for \cref{alg:pp} to
resolve (\cref{subsec:ablation}).

\noindent\textbf{Cost.}
For eligible robot $r$ and pool task $\tau$ the cost is
\begin{equation}\label{eq:cost}
  c(r,\tau) \;=\; \hat a(r,\tau)\,\bigl(1 + w_{\mathrm{b}}\,(1-s^r_t)\bigr),
\qquad
  \hat a(r,\tau) =
  \begin{cases}
    d(p^r,p_\tau) & \text{static (greedy, Hungarian)}\\[2pt]
    d(p^r,p_\tau) + w_{\mathrm{c}}\,\chi_{\mathcal{B}}(p^r\!\leadsto\!p_\tau) & \text{congestion proxy}\\[2pt]
    A^{\mathcal{B}}(p^r,p_\tau,t_0) & \text{path-aware (\sname)}
  \end{cases}
\end{equation}
with $c(r,\tau)=\infty$ when \cref{eq:feasible} fails. The battery factor
makes a low-charge robot look farther, so long trips go to well-charged
robots. The congestion proxy counts, along the free-space shortest path, the
cells reserved by other robots at the time the robot would traverse them
(held cells count $w_{\mathrm{h}}$); the path-aware cost is the arrival time
\cref{eq:arrival} of \cref{alg:astar} against the table that already contains
group~1. Costs are computed for the $K$ nearest pool tasks of each robot
(static distance), other pairs are infeasible.

\noindent\textbf{Assignment.}
The cost matrix is solved by the Hungarian method
\citep{kuhn1955hungarian,munkres1957algorithms} in the
shortest-augmenting-path form of \citet{jonker1987shortest}, applied to the
rectangular matrix, with infeasible
pairs priced above the sum of all finite costs; or, for the greedy baseline,
by repeatedly committing the cheapest remaining pair.

\begin{proposition}[Optimality of the assignment step]\label{prop:hungarian}
For a cost matrix $C\in(\mathbb{R}_{\ge0}\cup\{\infty\})^{m\times k}$, the
Hungarian step returns a matching that, among all matchings of maximum
cardinality over finite entries, minimises total cost; it runs in
$O(\min(m,k)^2\max(m,k))$ time.
\end{proposition}
\begin{proof}
Assume $m\le k$ (otherwise the implementation solves the transpose). Replace
$\infty$ by any $\Lambda>S=\sum_{ij:C_{ij}<\infty}C_{ij}$ (the implementation
uses $S+10^6$); the algorithm returns a minimum-cost matching that covers
every row. Such a matching $M$ with $\ell(M)$ $\Lambda$-entries costs
$\ell(M)\Lambda+F(M)$ with finite part $0\le F(M)\le S<\Lambda$, so a
minimiser first minimises $\ell$ and then $F$. Let $\nu$ be the maximum
cardinality of a matching of finite entries. Any such matching extends to a
row-covering matching because $m\le k$, and the added pairs are all
$\Lambda$-entries (a finite one would give a larger finite matching), so
$\min\ell=m-\nu$ and every maximum finite matching $M'$ appears, extended,
with cost $(m-\nu)\Lambda+F(M')$. Dropping the $\Lambda$-pairs of the
minimiser therefore leaves a maximum-cardinality finite matching of minimum
finite cost. Each of the $m$ rows is inserted by one shortest-augmenting-path
phase with potentials \citep{jonker1987shortest} of $O(mk)$ time, which gives
$O(m^2k)$.
\end{proof}

\begin{algorithm}[t]
\caption{\sname\ epoch (rolling horizon)}\label{alg:pact}
\KwIn{tick $t_0$, robots $\mathcal{R}$, pending tasks, window $w$}
$\mathcal{B}\gets$ empty table on $[t_0,t_0+w]$; release finished chargers; idle robots off their bay become parking\;
write holds for static robots (charging, loading, unloading, depleted, idle at bay)\;
\textsc{ChargerAssignment} (\cref{alg:charge})\;
group~1 $\gets$ busy robots $\cup$ parking robots on a busy goal or held last epoch; \textsc{PP}(group~1, $\mathcal{B}$) (\cref{alg:pp})\;
$\mathcal{E}\gets$ idle/parking robots not in group~1 with $s^r_t\ge\theta_{\mathrm{lo}}$; remove their holds from $\mathcal{B}$\;
build pool with endpoint capacities $\kappa_{\mathrm{d}},\kappa_{\mathrm{p}}$; $C_{r\tau}\gets c(r,\tau)$ by \cref{eq:cost} for $K$ candidates per $r\in\mathcal{E}$\;
$M\gets$ \textsc{Hungarian}$(C)$; assign $(r,\tau)\in M$: goal of $r\gets p_\tau$\;
unassigned $r\in\mathcal{E}$: goal of $r\gets h^r$ (parking)\;
group~2 $\gets$ newly assigned $\cup$ parking robots not in group~1; \textsc{PP}(group~2, $\mathcal{B}$, external = group~1)\;
\end{algorithm}

\noindent\textbf{Complexity.}
Let $|V|$ be the number of free cells. One epoch costs
$O\bigl(|\mathcal{R}|^2\,|V|\,w\bigr)$ for planning (\cref{prop:pp}, each
call bounded by $|V|(w+1)$ expansions; in practice the cascade is rare and
the cost is linear in $|\mathcal{R}|$), plus $O(mK\,|V|\,w)$ for path-aware
pricing and $O(m^2\,|\mathcal{T}_{\mathrm{pool}}|)$ for the assignment. The
static and proxy variants replace the pricing term by $O(mK)$ table look-ups
(distances are cached), which is the entire runtime overhead of the coupling;
\cref{subsec:main} reports it as planner milliseconds per simulated tick.
